% GENERATED FILE. Edit canonical lessons, then run npm run generate:books.
% canonical-source-hash: fnv1a64:84cdc5e1cb3107e2
% canonical-lessons: ML-C492-alavu, ML-C492-gunamenma, ML-C492-ulppannam, ML-C492-yatharttha, ML-C492-mayam-cerkkuka, ML-R492-review-485-492

\chapter{Size and Quality}
\label{ch:ml-size-and-quality}
\hlchaptermodality{492}

\begin{chapteropening}
\textbf{By the end of this chapter:} \emph{I can say a measurement, quality, a product, genuine and to adulterate.}

Complete the last lesson of chapter 492: I can say a measurement, quality, a product, genuine and to adulterate.
\end{chapteropening}

\section[\texorpdfstring{\ml{അളവ്} (aḷavŭ)}{അളവ് (aḷavŭ)}]{\ml{അളവ്} (aḷavŭ) — a measurement; a size}

\label{lesson:ML-C492-alavu}



\begin{quote}
Before the new one: say the Malayalam for a salesman, then the Malayalam for to try on.
\end{quote}

\subsection*{You'll want to know: \ml{അളവ്}}
\textbf{\ml{അളവ്}} — \emph{aḷavŭ} — \textquotedblleft{}a measurement; a size\textquotedblright{}.

\begin{quote}
\ml{തയ്യൽക്കാരൻ} \ml{എന്റെ} \ml{അളവ്} \ml{എടുത്തു}.
\end{quote}

\begin{quote}
\emph{tayyalkkāran enṟe aḷavŭ eṭuttu.}
\end{quote}

\begin{quote}
The tailor took my measurements.
\end{quote}

The noun of \ml{അളക്കുക}, to measure. A shirt in your size is \ml{നിങ്ങളുടെ} \ml{അളവിലുള്ള} \ml{ഷർട്ട്}.

\begin{cousinweb}
\textbf{\ml{അളവ്}} (\emph{aḷavŭ}) is native Dravidian, the noun of \ml{അളക്കുക}, to measure (DEDR 295).
\end{cousinweb}

\begin{grammarlens}[title={What you've built}]
One more word for getting around.
\end{grammarlens}

\subsection*{Your turn}
\begin{itemize}
\raggedright
  \item \emph{Say it:} \emph{aḷavŭ}
  \item \emph{Say it:} \emph{aḷavŭ}, once more
  \item \emph{Say it:} \emph{tayyalkkāran enṟe aḷavŭ eṭuttu.} — The tailor took my measurements.
  \item \emph{Recall:} say the Malayalam for a salesman, then the Malayalam for to try on, then say \emph{aḷavŭ} again
\end{itemize}

\begin{tcolorbox}[breakable,colback=teal!4,colframe=teal!35!black,title={Before you move on}]
What does \ml{അളവ്} mean? (\textquotedblleft{}A measurement; a size\textquotedblright{}.) Say it once more.
\end{tcolorbox}
\section[\texorpdfstring{\ml{ഗുണമേന്മ} (guṇamēnma)}{ഗുണമേന്മ (guṇamēnma)}]{\ml{ഗുണമേന്മ} (guṇamēnma) — quality}

\label{lesson:ML-C492-gunamenma}



\begin{quote}
Before the new one: say the Malayalam for to try on, then the Malayalam for a measurement.
\end{quote}

\subsection*{You'll want to know: \ml{ഗുണമേന്മ}}
\textbf{\ml{ഗുണമേന്മ}} — \emph{guṇamēnma} — \textquotedblleft{}quality\textquotedblright{}.

\begin{quote}
\ml{ആ} \ml{കടയിലെ} \ml{സാധനങ്ങൾക്ക്} \ml{നല്ല} \ml{ഗുണമേന്മയുണ്ട്}.
\end{quote}

\begin{quote}
\emph{ā kaṭayile sādhanaṅṅaḷkkŭ nalla guṇamēnmayuṇṭŭ.}
\end{quote}

\begin{quote}
The goods in that shop are of good quality.
\end{quote}

\ml{ഗുണം}, a good quality, plus \ml{മേന്മ}, excellence.

\begin{cousinweb}
\textbf{\ml{ഗുണമേന്മ}} (\emph{guṇamēnma}) is a Malayalam compound built from \ml{ഗുണം}, quality, from Sanskrit \emph{guṇa} (CDIAL 4190), and the native word \ml{മേന്മ}, excellence (DEDR 5086).
\end{cousinweb}

\begin{grammarlens}[title={What you've built}]
One more word for getting around.
\end{grammarlens}

\subsection*{Your turn}
\begin{itemize}
\raggedright
  \item \emph{Say it:} \emph{guṇamēnma}
  \item \emph{Say it:} \emph{guṇamēnma}, once more
  \item \emph{Say it:} \emph{ā kaṭayile sādhanaṅṅaḷkkŭ nalla guṇamēnmayuṇṭŭ.} — The goods in that shop are of good quality.
  \item \emph{Recall:} say the Malayalam for to try on, then the Malayalam for a measurement, then say \emph{guṇamēnma} again
\end{itemize}

\begin{tcolorbox}[breakable,colback=teal!4,colframe=teal!35!black,title={Before you move on}]
What does \ml{ഗുണമേന്മ} mean? (\textquotedblleft{}Quality\textquotedblright{}.) Say it once more.
\end{tcolorbox}
\section[\texorpdfstring{\ml{ഉൽപ്പന്നം} (ulppannaṁ)}{ഉൽപ്പന്നം (ulppannaṁ)}]{\ml{ഉൽപ്പന്നം} (ulppannaṁ) — a product}

\label{lesson:ML-C492-ulppannam}



\begin{quote}
Before the new one: say the Malayalam for a measurement, then the Malayalam for quality.
\end{quote}

\subsection*{You'll want to know: \ml{ഉൽപ്പന്നം}}
\textbf{\ml{ഉൽപ്പന്നം}} — \emph{ulppannaṁ} — \textquotedblleft{}a product\textquotedblright{}.

\begin{quote}
\ml{ആ} \ml{കമ്പനിയുടെ} \ml{ഉൽപ്പന്നങ്ങൾ} \ml{എല്ലാ} \ml{കടകളിലും} \ml{കിട്ടും}.
\end{quote}

\begin{quote}
\emph{ā kampaniyuṭe ulppannaṅṅaḷ ellā kaṭakaḷiluṁ kiṭṭuṁ.}
\end{quote}

\begin{quote}
That company's products are available in every shop.
\end{quote}

Plural: \ml{ഉൽപ്പന്നങ്ങൾ}, \emph{ulppannaṅṅaḷ}.

A -\ml{ം} noun: the plural is \ml{ഉൽപ്പന്നങ്ങൾ}.

\begin{cousinweb}
\textbf{\ml{ഉൽപ്പന്നം}} (\emph{ulppannaṁ}) is borrowed from Sanskrit \emph{utpanna}, \textquotedblleft{}risen, produced\textquotedblright{} (Monier-Williams): what has been made.
\end{cousinweb}

\begin{grammarlens}[title={What you've built}]
One more word for getting around.
\end{grammarlens}

\subsection*{Your turn}
\begin{itemize}
\raggedright
  \item \emph{Say it:} \emph{ulppannaṁ}
  \item \emph{Say it:} \emph{ulppannaṁ}, once more
  \item \emph{Say it:} \emph{ā kampaniyuṭe ulppannaṅṅaḷ ellā kaṭakaḷiluṁ kiṭṭuṁ.} — That company's products are available in every shop.
  \item \emph{Recall:} say the Malayalam for a measurement, then the Malayalam for quality, then say \emph{ulppannaṁ} again
\end{itemize}

\begin{tcolorbox}[breakable,colback=teal!4,colframe=teal!35!black,title={Before you move on}]
What does \ml{ഉൽപ്പന്നം} mean? (\textquotedblleft{}A product\textquotedblright{}.) Say it once more.
\end{tcolorbox}
\section[\texorpdfstring{\ml{യഥാർത്ഥ} (yathārttha)}{യഥാർത്ഥ (yathārttha)}]{\ml{യഥാർത്ഥ} (yathārttha) — genuine; real}

\label{lesson:ML-C492-yatharttha}



\begin{quote}
Before the new one: say the Malayalam for quality, then the Malayalam for a product.
\end{quote}

\subsection*{You'll want to know: \ml{യഥാർത്ഥ}}
\textbf{\ml{യഥാർത്ഥ}} — \emph{yathārttha} — \textquotedblleft{}genuine; real\textquotedblright{}.

\begin{quote}
\ml{ഇത്} \ml{യഥാർത്ഥ} \ml{തുകൽ} \ml{ആണോ}?
\end{quote}

\begin{quote}
\emph{itŭ yathārttha tukal āṇō?}
\end{quote}

\begin{quote}
Is this real leather?
\end{quote}

An adjective before a noun. The adverb is \ml{യഥാർത്ഥത്തിൽ}, in fact.

\begin{cousinweb}
\textbf{\ml{യഥാർത്ഥ}} (\emph{yathārttha}) is borrowed from Sanskrit \emph{yathārtha}, \textquotedblleft{}as the thing is\textquotedblright{}, true to reality (Monier-Williams).
\end{cousinweb}

\begin{grammarlens}[title={What you've built}]
One more word for getting around.
\end{grammarlens}

\subsection*{Your turn}
\begin{itemize}
\raggedright
  \item \emph{Say it:} \emph{yathārttha}
  \item \emph{Say it:} \emph{yathārttha}, once more
  \item \emph{Say it:} \emph{itŭ yathārttha tukal āṇō?} — Is this real leather?
  \item \emph{Recall:} say the Malayalam for quality, then the Malayalam for a product, then say \emph{yathārttha} again
\end{itemize}

\begin{tcolorbox}[breakable,colback=teal!4,colframe=teal!35!black,title={Before you move on}]
What does \ml{യഥാർത്ഥ} mean? (\textquotedblleft{}Genuine; real\textquotedblright{}.) Say it once more.
\end{tcolorbox}
\section[\texorpdfstring{\ml{മായം} \ml{ചേർക്കുക} (māyaṁ cērkkuka)}{മായം ചേർക്കുക (māyaṁ cērkkuka)}]{\ml{മായം} \ml{ചേർക്കുക} (māyaṁ cērkkuka) — to adulterate}

\label{lesson:ML-C492-mayam-cerkkuka}



\begin{quote}
Before the new one: say the Malayalam for a product, then the Malayalam for genuine.
\end{quote}

\subsection*{You'll want to know: \ml{മായം} \ml{ചേർക്കുക}}
\textbf{\ml{മായം} \ml{ചേർക്കുക}} — \emph{māyaṁ cērkkuka} — \textquotedblleft{}to adulterate\textquotedblright{}.

\begin{quote}
\ml{ചില} \ml{കടക്കാർ} \ml{പാലിൽ} \ml{മായം} \ml{ചേർക്കുന്നു}.
\end{quote}

\begin{quote}
\emph{cila kaṭakkār pālil māyaṁ cērkkunnu.}
\end{quote}

\begin{quote}
Some shopkeepers adulterate the milk.
\end{quote}

Past: \ml{മായം} \ml{ചേർത്തു}, \emph{māyaṁ cērttu}.

\ml{മായം} is adulteration; \ml{ചേർക്കുക} is to add. What is spoilt takes -\ml{ഇൽ}: \ml{പാലിൽ}.

\begin{cousinweb}
\textbf{\ml{മായം} \ml{ചേർക്കുക}} (\emph{māyaṁ cērkkuka}) is a Malayalam verb phrase built from \ml{മായം}, adulteration, usually traced to Sanskrit \emph{māyā}, \textquotedblleft{}art, trickery\textquotedblright{} (Monier-Williams), and the native verb \ml{ചേർക്കുക}, to add (DEDR 2814).
\end{cousinweb}

\begin{grammarlens}[title={What you've built}]
One more word for getting around.
\end{grammarlens}

\subsection*{Your turn}
\begin{itemize}
\raggedright
  \item \emph{Say it:} \emph{māyaṁ cērkkuka}
  \item \emph{Say it:} \emph{māyaṁ cērkkuka}, once more
  \item \emph{Say it:} \emph{cila kaṭakkār pālil māyaṁ cērkkunnu.} — Some shopkeepers adulterate the milk.
  \item \emph{Recall:} say the Malayalam for a product, then the Malayalam for genuine, then say \emph{māyaṁ cērkkuka} again
\end{itemize}

\begin{tcolorbox}[breakable,colback=teal!4,colframe=teal!35!black,title={Before you move on}]
What does \ml{മായം} \ml{ചേർക്കുക} mean? (\textquotedblleft{}To adulterate\textquotedblright{}.) Say it once more.
\end{tcolorbox}
\section[(dialogue)]{Words from chapters 485–492 — a review}

\label{lesson:ML-R492-review-485-492}



\begin{quote}
Say the Malayalam for a lift and a staircase.
\end{quote}

\subsection*{Your turn}
Say these aloud:

\begin{itemize}
\raggedright
  \item a lift, a staircase, a gateway, a waiting room, a roadside food stall
  \item an embassy, a department, a branch office, office hours, a working day
  \item a municipality, a village office, a taluk, a district collector, a nuisance
  \item a signature, a duplicate copy, an identity document, a validity period, to expire
  \item currency, exchange, to break a note, to make a payment, to stand in line
  \item a post box, a courier, to receive officially, to send back, to contact
  \item a shopping mall, a special offer, a clearance sale, a salesman, to try on
  \item a measurement, quality, a product, genuine, to adulterate
\end{itemize}

\begin{tcolorbox}[breakable,colback=teal!4,colframe=teal!35!black,title={Before you move on}]
A lift? (\textbf{\ml{ലിഫ്റ്റ്}}, \emph{liphṟṟŭ}.) A staircase? (\textbf{\ml{കോണിപ്പടി}}, \emph{kōṇippaṭi}.) A gateway? (\textbf{\ml{കവാടം}}, \emph{kavāṭaṁ}.) A waiting room? (\textbf{\ml{വിശ്രമമുറി}}, \emph{viśramamuṟi}.) A roadside food stall? (\textbf{\ml{തട്ടുകട}}, \emph{taṭṭukaṭa}.) An embassy? (\textbf{\ml{എംബസി}}, \emph{embasi}.) A department? (\textbf{\ml{വകുപ്പ്}}, \emph{vakuppŭ}.) A branch office? (\textbf{\ml{ശാഖ}}, \emph{śākha}.) Office hours? (\textbf{\ml{പ്രവൃത്തിസമയം}}, \emph{pravṛttisamayaṁ}.) A working day? (\textbf{\ml{പ്രവൃത്തിദിനം}}, \emph{pravṛttidinaṁ}.) A municipality? (\textbf{\ml{നഗരസഭ}}, \emph{nagarasabha}.) A village office? (\textbf{\ml{വില്ലേജ്} \ml{ഓഫീസ്}}, \emph{villējŭ ōphīsŭ}.) A taluk? (\textbf{\ml{താലൂക്ക്}}, \emph{tālūkkŭ}.) A district collector? (\textbf{\ml{കളക്ടർ}}, \emph{kaḷakṭar}.) A nuisance? (\textbf{\ml{ശല്യം}}, \emph{śalyaṁ}.) A signature? (\textbf{\ml{ഒപ്പ്}}, \emph{oppŭ}.) A duplicate copy? (\textbf{\ml{പകർപ്പ്}}, \emph{pakarppŭ}.) An identity document? (\textbf{\ml{തിരിച്ചറിയൽരേഖ}}, \emph{tiriccaṟiyalrēkha}.) A validity period? (\textbf{\ml{കാലാവധി}}, \emph{kālāvadhi}.) To expire? (\textbf{\ml{കാലഹരണപ്പെടുക}}, \emph{kālaharaṇappeṭuka}.) Currency? (\textbf{\ml{കറൻസി}}, \emph{kaṟansi}.) Exchange? (\textbf{\ml{വിനിമയം}}, \emph{vinimayaṁ}.) To break a note? (\textbf{\ml{ചില്ലറയാക്കുക}}, \emph{cillaṟayākkuka}.) To make a payment? (\textbf{\ml{പണമടയ്ക്കുക}}, \emph{paṇamaṭaykkuka}.) To stand in line? (\textbf{\ml{വരിനിൽക്കുക}}, \emph{varinilkkuka}.) A post box? (\textbf{\ml{തപാൽപ്പെട്ടി}}, \emph{tapālppeṭṭi}.) A courier? (\textbf{\ml{കൊറിയർ}}, \emph{koṟiyar}.) To receive officially? (\textbf{\ml{കൈപ്പറ്റുക}}, \emph{kaippaṟṟuka}.) To send back? (\textbf{\ml{തിരിച്ചയയ്ക്കുക}}, \emph{tiriccayaykkuka}.) To contact? (\textbf{\ml{ബന്ധപ്പെടുക}}, \emph{bandhappeṭuka}.) A shopping mall? (\textbf{\ml{മാൾ}}, \emph{māḷ}.) A special offer? (\textbf{\ml{ഓഫർ}}, \emph{ōphar}.) A clearance sale? (\textbf{\ml{ആദായവിൽപ്പന}}, \emph{ādāyavilppana}.) A salesman? (\textbf{\ml{വിൽപ്പനക്കാരൻ}}, \emph{vilppanakkāran}.) To try on? (\textbf{\ml{ഇട്ടുനോക്കുക}}, \emph{iṭṭunōkkuka}.) A measurement? (\textbf{\ml{അളവ്}}, \emph{aḷavŭ}.) Quality? (\textbf{\ml{ഗുണമേന്മ}}, \emph{guṇamēnma}.) A product? (\textbf{\ml{ഉൽപ്പന്നം}}, \emph{ulppannaṁ}.) Genuine? (\textbf{\ml{യഥാർത്ഥ}}, \emph{yathārttha}.) To adulterate? (\textbf{\ml{മായം} \ml{ചേർക്കുക}}, \emph{māyaṁ cērkkuka}.)
\end{tcolorbox}
